\documentstyle[%


righttag%


,11pt%


,amssymb%


,russian%               remove in Eng.


,izvru%                 Set izven% in Eng.


]{amsart}





% %repl. 0,5 for 0.5





\newcommand{\const}{\operatorname{const}}


\newcommand{\tts}[1]{}





%\hoffset=-15mm%                  For Eng.


 


\setcounter{page}{41}


\god{1999}


\nomer{11~(450)} %                        Remove in Eng.


%\nomer{Vol.\,43, No.\,11}%               Set in Eng.


%\str{\pageref{begin}--\pageref{end}}%  Set in Eng.


\udk{517.982}





\author{\it В.Л.\,Крепкогорский}


% NB:   ^^ remove in Eng.


\title{О многомерных методах интерполяции}





\date{}


%\pagestyle{myheadings}%         Set in Eng.


\pagestyle{plain} %              Remove in Eng.


\begin{document}


%\thispagestyle{plain}%          Set in Eng.


\maketitle


\label{begin}





Известен ряд работ, в которых авторы предлагают вещественные методы,


аналогичные обычным, но с функторами, действующими не на


пару пространств, а на набор из $n$ или $2^n$ пространств. Будем называть


такие методы многомерными. Многомерные аналоги вещественного


$K$-метода рассматривались в [1]--[4].


Заметим, что целесообразность введения многомерных методов


до сих пор не продемонстрирована достаточно убедительно, т.\,к. эти методы


реализованы только в очень редких случаях. В данной работе рассматривается


простейший многомерный метод. Его функторы ${\cal F}_{\Theta,Q}$ составлены


из обычных функторов $K$-метода. В классах пространств с двумя


параметрами этот метод удобно реализовать, интерполируя сначала по


одному параметру, а затем --- по другому. Например, при интерполяции в


классе пространств Бесова будем сначала интерполировать по


параметру $p$, а затем --- по $s$. При этом на первом этапе получаются


пространства, которые не принадлежат даже расширенному классу


пространств Бесова $B^s_{p,q,(r)}$. Однакo на следующем этапе


снова возвращаемся в класс пространств Бесова. Этот класс оказывается


стабильным относительно функторов ${\cal F}_{\Theta, Q}$, хотя при


одномерной интерполяции с изменением параметра $p$ мы, вообще говоря,


покидаем класс пространств Бесова.





С помощью многомерных функторов получена интерполяционная теорема


для пространств $B^s_p$, использующая слабые условия вида


\begin{equation}


T:B^s_{p,1,(1)}\rightarrow B{}^{\wt s}_{\wt p,\infty,(\infty) }.


\end{equation}


Заметим, что одномерная теорема для пространств $B^s_p$ со


слабыми условиями вида (1), вообще говоря, неверна. Известен


пример оператора, который удовлетворяет слабым условиям при любых


$(p,s)$, $(\wt p, \wt s)$, для которых выполняются условия


$-\infty < s< +\infty$, $1<p<\infty$, $s=\wt s =1/p$, $p=\wt p$, но


$T\notin {\cal L}(B^{1/r}_r,B^{1/r}_r)$ ни при каком $r\in (1,\infty)$.





Отметим, наконец, что в [5] содержится утверждение


о якобы имеющем место совпадении методов ${\cal F}_{\Theta, Q}$ и


Фернандеса. Однако целый ряд авторов [6], [7], [8], [4] привели примеры,


показывающие что, вообще говоря, это не так.





\section{Основные определения и обозначения}





Пусть даны четыре банаховых пространства $A_{0,0}$, $A_{1,0}$,


$A_{0,1}$, $A_{1,1}$, вложенных в одно хаусдорфово ТВП $\cal A$. На


четверке пространств $A_{i,k}$ определим многомерный функтор


$$


{\cal F}_{\Theta, Q}((A_{i,k})_{k=0,\,1})_{i=0,\,1}:=


((A_{0,0}, A_{0,1})_{\theta_1, q_1}, (A_{1,0}, A_{1,1})_{\theta_1,


q_1})_{\theta_2, q_2};


$$


$\Theta =(\theta_1, \theta_2)$, $Q=(q_1, q_2)$, $0< \theta_1,


\theta_2 <1$, $1< q_1,q_2 <\infty$.





Назовем квазинормированной (банаховой) решеткой


квазинормированное (банахово) пространство функций $E$, норма которого


обладает свойством монотонности


$$


|f(x)| \le |g(x)|\ \; \forall x \Rightarrow \|f\mid E\| \le \|g\mid E\|.


$$





Говорят, что банахова решетка $E$ обладает свойством Фату, если для


монотонно возрастающей последовательности $f_n \uparrow f$ из условия


$\lim\limits_{n\rightarrow \infty }\|f_n\mid E\| <\infty$ следует $f\in E$ и


$\|f\mid E\|=\linebreak


\lim\limits_{n\rightarrow \infty }\|f_n\mid E\| <\infty$.





Пусть даны две квазинормированные решетки $E(X)$ и $G(Y)$, а


также функция двух переменных $f(x,y)$, $x\in X$, $y\in Y$. Смешанной


квазинормой назовем функционал $\|f\mid G[E]\|=\linebreak


\big\| \|f(x,y) \mid E(X)\|\mid G(Y)\big\|$.





Обозначим через $L_{p,q}$


пространство Лоренца, через $B^s_{p,q,(r)}(R_n)$ --- пространство


Бесова и $F^s_{p,q,(r)}(R_n)$ --- пространство Лизоркина--Трибеля с


нормами


\begin{align*}


\|(a_j)^{\infty}_{j=0}\mid\ell^s_q\| &:= \| a_j\cdot 2^{js}\mid\ell_q\|;\\


%


\|f\mid L_{p,q}\| &:= \bigg( \int_0^{\infty} (t^{1/p}


f^*(t))^q\frac{dt}{t} \bigg)^{1/q},


\end{align*}


где $f^*$ --- невозрастающая равноизмеримая перестановка функции $f$;


\begin{align*}


\| f\mid B^s_{p,q,(r)}\| &:=


\| F^{-1}\varphi_j Ff\mid\ell^s_q[L_{p,r}]\|\quad (\varphi_j)\in \Phi ;\\


%


\| f\mid F^s_{p,q,(r)}\| &:= \| F^{-1}\varphi_j Ff\mid L_{p,r}[\ell^s_q]\|


\quad (\varphi_j)\in \Phi.


\end{align*}


Здесь $\Phi$ --- специальный класс последовательностей функций ([8],


c.\,203).


Для функций двух переменных $f(x_1,x_2) $ через $f^*_1$ ($f^*_2$)


обозначим равноизмеримую невозрастающую перестановку функции одной


переменной, которая получается если фиксировать переменную $x_2$ ($x_1$).





Если функция $f(x)$ определена на множестве, на котором


рассматриваются две различные меры $\mu$ и $\lambda$, то для


определенности будем писать $f^{*(\mu )}$ в случае, когда


перестановка сделана равноизмеримой относительно меры $\mu$.





При интерполяции пространств Бесова получаются пространства


$BL^{s,k}_{p,q}$ [9]


$$


(B^{s_0}_{p_0},B^{s_1}_{p_1})_{\theta,q}=


BL^{s,k}_{p,q} =\{ f\in S':\| f\mid BL\| = \| (f*\varphi_j)


\mid L^{s,k}_{p,q}\| <\infty \},


$$


где $L^{s,k}_{p,q} = L_{p,q}(2^{j(s-k/p)}$, $m_n\times (2^{jk}\nu))$


--- пространство Лоренца $L_{p,q}$ c весом $2^{j(s-k/p)}$, $m_n$ ---


лебегова мера на $R_n$, $2^{jk}\nu \{ j\}$ --- значение атомической


меры, и $-\infty <s<\infty$, $-\infty <k<\infty$, $1<p<\infty$,


$0<q\le \infty$, $1/p=(1-\theta )/p_0+


\theta /p_1$, $s=(1-\theta )s_0+\theta s_1$, $k$ --- угловой коэффициент


прямой, проходящей через точки $(1/p_i,s_i)$, $i=0,\,1$.





\section{Реализация метода ${\cal F}_{\Theta, Q} $ в классах


пространств Бесова\protect\\ и Лизоркина--Трибеля }





Cформулируем в виде леммы два простых утверждения.





\begin{lem}


Пусть $E$ --- банахова структура с нормой, обладающей


свойством Фату, тогда


\begin{alignat*}{2}


&{\mathrm a)} &\ \big\| \,\| f_n(y)\mid\ell_1(N) \| \mid E(Y)\big\|


&\le\big \| \, \|f_n(y)\mid E(Y)\| \mid \ell_1(N)\big\| ; \\


%


&{\mathrm b)} &\ \big\| \,\| f_n(y)\mid\ell_{\infty} (N)


\| \mid E(Y)\big\| &\ge \big\| \, \|f_n(y)


\mid E(Y)\| \mid \ell_{\infty } (N)\big\|.


\end{alignat*}


\end{lem}





Пространства $\ell^s_q[L_{p,r}] $ хорошо интерполируются по


``внешним'' параметрам $s$ и $q$. Опишем, что получается при интерполяции


пространств со смешанной нормой по ``внутренним'' параметрам.





\begin{lem}


Пусть $\ell$ --- банахова структура


последовательностей и $(a_j)$, $a_j \in R^1$, $1\le r_0,r_1\le \infty$,


$1<p_0,p_1< \infty$, тогда


$$


K(t,(\psi_j),\ell [L_{p_0,r_0}], \ell[L_{p_1,r_1}]) \sim \|


K(t,(\psi_j), L_{p_0,r_0}, L_{p_1,r_1})\mid\ell \|.


$$


\end{lem}





{\bf Доказательство.} Оценим $K$-функционал


\begin{multline*}


K(t,(\psi_j),\ell [L_{p_0,r_0}], \ell[L_{p_1,r_1}]) =


\inf_{\psi^0_j +\psi^1_j =\psi_j}(\big\|\,


\|\psi^0_j\mid L_{p_0,r_0}\|\mid\ell \big\| + \big\| \, \| \psi^1_j


\mid L_{p_1,r_1}\| \mid\ell \big\| ) \ge \\


%


\ge \inf_{\psi^0_j +\psi^1_j =\psi_j}(\|\,


( \|\psi^0_j\mid L_{p_0,r_0}\| + \| \psi^1_j


\mid L_{p_1,r_1}\|)\mid\ell \| ) \ge


\| \inf_{\psi^0_j +\psi^1_j =\psi_j}(


\|\psi^0_j\mid L_{p_0,r_0}\| + \| \psi^1_j


\mid L_{p_1,r_1}\|) \mid\ell \| = \\


%


= \| K(t,(\psi_j), L_{p_0,r_0},L_{p_1,r_1} ) \mid \ell \|.


\end{multline*}


С другой стороны, выбирая при фиксированном $j$ пару $\psi^0_j$ и


$\psi^1_j$ таким образом, чтобы выполнялись условия


$\psi^0_j+\psi^1_j =\psi_j $ и


$$


\| \psi^0_j\mid L_{p_0,r_0}\| +t\| \psi^1_j \mid L_{p_1,r_1}\| \le


2K(t,(\psi_j),L_{p_0,r_0},L_{p_1,r_1}),


$$


получим оценку


$$


\|\,( \| \psi^0_j\mid L_{p_0,r_0}\| +t\| \psi^1_j \mid L_{p_1,r_1}\|)\mid


\ell \| \le 2\| K(t,(\psi_j),L_{p_0,r_0},L_{p_1,r_1})\mid \ell \|,


$$


откуда следует


\begin{multline*}


\big\| \, \| \psi^0_j\mid L_{p_0,r_0}\| \mid \ell \big\|


+t\big\| \, \| \psi^1_j \mid L_{p_1,r_1}\|\mid\ell \big\| \le


2 \| \max ( \| \psi^0_j\mid L_{p_0,r_0}\|;\ t\| \psi^1_j


\mid L_{p_1,r_1}\| ) \mid \ell \| \le \\


%


\le 2 \| ( \| \psi^0_j\mid L_{p_0,r_0}\| + t\| \psi^1_j


\mid L_{p_1,r_1}\| ) \mid \ell \| \le


4\| K(t,(\psi_j),L_{p_0,r_0},L_{p_1,r_1})\mid \ell \|.\quad\square


\end{multline*}





Применим функтор ${\cal F}_{\Theta, Q}$ к


пространствам $\ell^s_q[L_{p,r}]$.





\begin{lem}


Пусть $\Theta =(\theta_1, \theta_2)$, $Q=(q_1,q_2)$,


$0<\theta_1, \theta_2<1$, $1\le q_1,q_2\le \infty$, $-\infty <s_k<\infty$,


$1<p<\infty$, $i=0,\,1;$ $k=0,\,1$. Если


$1/p=(1-\theta_1)/p_0+\theta_1/p_1$, $s=(1-\theta_2)+\theta_2s_1$, то


$$


{\cal F}_{\Theta,


Q}((\ell^{s_k}_{q_{i,k}}[L_{p_i,r_{i,k}}])_{i=0,\,1,\; k=\const})_{k=0,\,


1}= \ell^s_{q_2}[L_{p,q_1}].


$$


\end{lem}





{\bf Доказательство.} Как известно, $\ell^s_1 \subset \ell^s_q


\subset \ell^s_{\infty }$ и $L_{p,1} \subset L_{p,r} \subset L_{p,


\infty}$, поэтому $\ell^s_1[L_{p,1}]\subset \ell^s_q[L_{p,r}]


\subset \ell^s_{\infty}[L_{p, \infty }]$. Следовательно,


достаточно показать, что


$$


{\cal F}_{\Theta,Q}((\ell^{s_k}_1[L_{p_i,1}])_{i=0,\,1})_{k=0,\,1}=


{\cal F}_{\Theta,Q}((\ell^{s_k}_{\infty


}[L_{p_i,\infty}])_{i=0,\,1})_{k=0,\,1}=


\ell^s_{q_2}[L_{p,q_1}].


$$


Положим $u=1$ или $\infty$. Тогда с помощью леммы 2 можно


описать норму, получающуюся на первом этапе


$$


\| (\psi_j)\mid(\ell^{s_k}_u[L_{p_0,u}],


\ \ell^{s_k}_u[L_{p_1,u}])_{\theta_1,q_1} \| \sim


\Phi_{\theta_1,q_1}(\|K(t,\psi_j, L_{p_0,u}, L_{p_1,u})\mid\ell^{s_k}_u\|).


$$


Обозначим кратко структуру с такой нормой через $E_{s_k}$.


По лемме 1


\begin{multline*}


\|(\psi_j)\mid\ell^{s_k}_{\infty }[L_{p,q}]\| =


\| \Phi_{\theta_1,q_1}(K(t,(\psi_j), L_{p_0,u}, L_{p_1,u}))


\mid\ell^{s_k}_{\infty }\| \le


C\Phi_{\theta_1,q_1}(\| K(t,(\psi_j), L_{p_0,u}, L_{p_1,u})


|\ell^{s_k}_{\infty }\| )\le \\


%


\le C\| (\psi_k)\mid E_{s_k} \| \le C\Phi_{\theta_1,q_1}(\| K(t,(\psi_j),


L_{p_0,u}, L_{p_1,u})\mid \ell^{s_k}_1\| ) \le \\


%


\le C\| \Phi_{\theta_1,q_1}( K(t,(\psi_j), L_{p_0,u},


L_{p_1,u}))\mid \ell^{s_k}_1\| \sim


C\|(\psi_j)\mid\ell^{s_k}_1[L_{p,q}]\|,


\end{multline*}


т.\,е. $\ell^{s_k}_1[L_{p,q}] \subset E_{s_k} \subset


\ell^{s_k}_{\infty}[L_{p,q}]$. Поэтому из равенства


$$


(\ell^{s_0}_1[L_{p,q}], \ell^{s_1}_1[L_{p,q}])_{\theta_2,q_2} =


(\ell^{s_0}_{\infty }[L_{p,q}],


\ell^{s_1}_{\infty }[L_{p,q}])_{\theta_2,q_2}=


\ell^s_{q_2}[L_{p,q_1}]


$$


следует $(E_{s_0},E_{s_1})_{\theta_2,q_2}=\ell^s_{q_2}[L_{p,q}] $ и


\begin{multline*}


{\cal F}((\ell^{s_k}_{q_{i,k}}[L_{p_i,r_{i,k}}])_{i=0,\,1})_{k=0,\,1}=\\


%


((\ell^{s_0}_u[L_{p_0,u}], \ell^{s_0}_u[L_{p_1,u}])_{\theta_1,q_1},


(\ell^{s_1}_u[L_{p_0,u}],


\ell^{s_1}_u[L_{p_1,u}])_{\theta_1,q_1})_{\theta_2,q_2}=


(E_{s_0},E_{s_1})_{\theta_2,q_2}=\ell^s_{q_2}[L_{p,q_1}].


\end{multline*}





Сформулируем теперь интерполяционную теорему ${\cal F}_{\Theta,Q}$--метода


для пространств Бесова.





\begin{thm}


Пусть $\Theta = (\theta_1, \theta_2)$, $Q=(q_1,q_2)$,


$0<\theta_1, \theta_2<1$, $1\le q_1,q_2\le \infty$,


$-\infty <s_k<\infty$, $1\le q_{i,k} \le \infty$, $1<p<\infty$,


$i=0,\, 1;$ $k=0, \,1$.


Если $1/p= (1-\theta_1)/p_0+\theta_1/p_1$,


$s=(1-\theta_2)s_0+\theta_2s_1$, то


$$


{\cal F}_{\Theta,Q}((B^{s_k}_{p_i,q_{i,k},(r_{i,k})})_{i=0,\,1,


\;k=\const})_{k=0,\,1} = B^s_{p,q_2,(q_1)}.


$$


\end{thm}





Доказательство этого утверждения следует из леммы 3 и теоремы об


интерполяции ретракций.





По лемме 2 на первом этапе применения функторов


$(\cdot, \cdot)_{\theta, q}$


получаем интерполяционные пространства с нормами вида


$$


\| f\|=\Phi_{\theta_1,q_1}(\| K(t, (\psi_j),


L_{p_0,u},L_{p_1,u})\mid\ell^{s_k}_1\| ).


$$


Такие пространства уже не принадлежат классу пространств


$B^s_{p,q,(r)} $ и только после повторного применения функторов


$(\cdot,\cdot )_{\theta,q}$ снова возвращаемся в класс пространств


Бесова.





Таким образом, при покоординатной интерполяции класс пространств


Бесова оказывается стабильным относительно функторов ${\cal F}_{\Theta,Q}$.


В то же время класс пространств Бесова, вообще говоря,


нестабилен относительно функторов $(\cdot,\cdot )_{\theta,q}$.


Пространства Бесова при одномерной интерполяции получаем


только в частных случаях а)~$(B^{s_0}_p, B^{s_1}_p)_{\theta,q}$ и


б)~$(B^{s_0}_{p_0}, B^{s_1}_{p_1})_{\theta,\wt q}$ при единственном


значении параметра $\wt q$, при котором


$1/\wt q= (1-\theta)/p_0 +\theta /p_1$.


В общем случае при одномерной интерполяции получаем


пространства типа $BL^{s,k}_{p,q}$. Теорема 1 позволяет получить


интерполяционную теорему со слабыми условиями типа (1)


$$


T:B^s_{p,1,(1)}\rightarrow B{}^{\wt s}_{\wt p,\infty,(\infty)}.


$$


Очевидно, можно рассмотреть другие типы слабых условий, например,


$T:F^s_{p,1}\rightarrow F{}^{\wt s}_{\wt p,\infty}$ или


$T:BL^{s,k}_{p,1} \rightarrow BL{}^{\wt s, \wt k}_{\wt p,\infty}$.





Следующая теорема показывает, что условия (1) самые


слабые, которые можно сформулировать в терминах банаховых пространств


типов $B$, $F$, $BL$.





\begin{thm}


При любых $1<p, \wt p <\infty$, $-\infty <s,\wt s,k<\infty$


имеют место вложения


\begin{alignat*}{4}


&\text{\em а)}&\ \; B^s_{p,1,(1)} &\subset BL^{s,k}_{p,1};


&\quad &\text{\em б)}&\ \;BL^{s,k}_{p,\infty}


&\subset B^s_{p,\infty,(\infty )}; \\


%


&\text{\em в)}&\ \; B^s_{p,1,(1)}&\subset F^s_{p,1};


&\quad &\text{\em г)}&\ \; F^s_{p,\infty }&\subset


B^s_{p,\infty,(\infty)}.


\end{alignat*}


\end{thm}





\begin{pf}


Вложения б) и г) следуют из очевидных вложений


\begin{alignat*}{2}


&\text{б}') &\ \ell^s_1[L_{p,1}] &\subset \ell^s_1[L_p]


\subset L_p[\ell^s_1]; \\


%


&\text{г}')&\ L_p[\ell^s_{\infty} ] &\subset \ell^s_{\infty }[L_p] \subset


\ell^s_{\infty }[L_{p,\infty }].


\end{alignat*}


Для проверки вложений а) и в) достаточно доказать вложения


$$


\text{а}')\ \ell^s_1[L_{p,1}] \subset L^{s,k}_{p,1} \quad \text{и} \quad


\text{б}')


\ L^{s,k}_{p,\infty } \subset \ell^s_{\infty }[L_{p,\infty }].


$$


Проверим справедливость a$'$). Пространство $\ell^s_1[L_{p,1}]$


состоит из последовательностей функций $\{ \psi_j\}^{\infty}_{j=0}$.


Любую положительную, измеримую функцию $\psi$


можно представить в виде предела монотонно возрастающей


последовательности простых функций $\xi_k=\sum\limits^m_{i-1}


C_{i,k}\chi_{_{H_{i,k}}}(x)$, $C_{i,k}>0$, $C_{i,k}=\const$,


$x\in R_n$. Тогда $\xi_k\uparrow \psi$ и


$$


\| \psi \mid L_{p,1} \|=\lim_{k\rightarrow \infty }\| \xi_k\mid L_{p,1}\|=


\lim_{k\rightarrow \infty }\bigg\| \sum^{m(k)}_{i=1}


C_{i,k}\chi_{_{H_{ik}}}(x)\mid L_{p,1} \bigg\|=


\sum^{m(k)}_{i=1} C_{i,k}\| \chi_{_{H_{ik}}}(x)\mid L_{p,1} \|.


$$


Можно считать, что последовательность


$\{\psi_j\}^{\infty}_{j=0} \in \ell^s_1[L_{p,1}]$


состоит из положительных функций.


Тогда существует последовательность простых функций


$\xi_{ik}\uparrow\psi_j $ при $k\rightarrow \infty$. В этом случае


\begin{multline*}


\|(\psi_j) \mid \ell^s_1[L_{p,1}] \| = \lim_{k\rightarrow \infty}\|


\xi_{jk} \mid\ell^s_1[L_{p,1}]\| =\\


%


=\lim_{k\rightarrow \infty }\bigg\| \sum^{m(j,k)}_{i=1}


C_{i,j,k}\chi_{_{H_{i,j,k}}}(x)\mid\ell^s_1[L_{p,1}] \bigg\|=


\lim_{k\rightarrow \infty } \sum^{m(j,k)}_{i=1}


C_{i,j,k}\| \chi_{_{H_{i,j,k}}}(x)\mid\ell^s_1[L_{p,1}] \|.


\end{multline*}


Оценим нормы характеристических функций


\begin{multline*}


\| \chi_{(j)\times H}(j,x)\mid L^{s,k}_{p,1} \| = \| 2^{j(s-k/p)}


\chi_{[0,\mu (H)\cdot 2^{jk}]} (t)\mid L_{p,1}\| = \\


%


=M(p) \mu^{1/p}(H) 2^{jk/p} 2^{j(s-k/p)} =


M(p)\,2^{js}\mu^{1/p}(H)


\end{multline*}


и


$$


\| \chi_{(j)\times H}\mid\ell^s_1[L_{p,1}]\| = M(p)


\|\mu^{1/p}(H) \chi_{(j)} \mid\ell^s_1 \| =


M(p) 2^{js}\mu^{1/p}(H).


$$


Здесь $M(p) =\| \chi_{[0,a]}\mid L_{p,1}\| /a^{1/p} =p^2/(p-1)$.


Значит, $\| \chi_{(j)\times H}\mid\ell^s_1[L_{p,1}]\| =


\|\chi_{(j)\times H}\mid L_{p,1}^{s,k}\|$. Тогда можно сделать окончательную


оценку


\begin{multline*}


\| (\psi_j)|L^{s,k}_{p,1}\| = \lim_{k\rightarrow \infty }


\bigg\| \sum^{m(j,k)}_{i=1} C_{i,j,k}\chi_{_{H_{i,j,k}}}(x)\mid


L_{p,1}^{s,k} \bigg\| \le \\


%


\le\lim_{k\rightarrow \infty } \sum^{m(j,k)}_{i=1}


C_{i,j,k}\| \chi_{_{H_{i,j,k}}}(x)\mid L^{s,k}_{p,1} \|=


\lim_{k\rightarrow \infty } \sum^{m(j,k)}_{i=1}


C_{i,j,k}\| \chi_{_{H_{i,j,k}}}(x)\mid\ell^s_1[L_{p,1}] \|=


\| (\psi_j) \mid\ell^s_1[L_{p,1}]\|.


\end{multline*}


Таким образом, вложение a$')$ доказано.





Для банаховой структуры $E$ через $E^1$ обозначим ассоциированное


пространство всех линейных функционалов, имеющих интегральное


представление. Очевидно, что из вложения $E\subset G$ следует


$G^1\subset E^1$. Положим $1/p +1/p'=1$. Известно [9], что


$(L^{-s,k}_{p',1})^1 = L^{s,k}_{p,\infty}$ и


$(\ell^{-s}_1[L_{p',1}])^1=\ell^s_{\infty }[L_{p,\infty }]$. Поэтому из


вложения $\ell^{-s}_1[L_{p',1}]\subset L^{-s,k}_{p',1}$ следует


$L^{s,k}_{p,\infty }\subset \ell^s_{\infty} [L_{p,\infty }]$.


\end{pf}





Пусть преобразование $A:R_2\rightarrow R_2$ определено формулой


$$


A(x,y) :=


\begin{pmatrix} a_1 & 0\\0 & a_2 \end{pmatrix}


\begin{pmatrix} x \\ y \end{pmatrix} +


\begin{pmatrix} h_1 \\ h_2 \end{pmatrix}, \quad


a_1,a_2,h_1,h_2 \in R_1.


 $$


Обозначим координаты образа $A(1/p,\,s)$ через $\wt \alpha$ и


$\wt s$. Пусть также $\wt p =1/\wt \alpha$. Преобразование $A$


переводит прямые в прямые, сохраняет соотношение отрезков и


параллельность осям координат.





Сформулируем теорему, связывающую сильные и слабые условия для


точек $(1/p,s)$ и $A(1/p,s)=(1/\wt p,\wt s)$.





\begin{thm}


Пусть $\Theta = (\theta_1,\theta_2)$,


$Q=(q_1,q_2)$, $0<\theta_1,\theta_2<1$, $1\le q_1,q_2\le \infty$,


$1<p_i,\wt p_i <\infty$, $-\infty <s_j,\wt s_j<\infty$,


$i=0,\,1$,j$=0,\,1;$


\begin{equation}


1/p=(1-\theta_1)/p_0+\theta_1/p_1,\quad


s=(1-\theta_2)s_0+\theta_2s_1.


\end{equation}


Если линейный оператор


\begin{equation}


T:B^{s_j}_{p_i,1,(1)} \rightarrow B{}^{\wt s_j}_{\wt


p_i,\infty,(\infty )},


\end{equation}


то $T:B^s_{p,q,(r)} \rightarrow B{}^{\wt s}_{\wt p,q,(r)}$.


Если, кроме того, $p\le \wt p$, то


$T:B^s_p \rightarrow B{}^{\wt s}_{\wt p}$.


\end{thm}








{\bf Доказательство.} Заметим, что условия (2) выполняются не только


для $(1/p_i,s_j)$, но и для параметров $ (1/\wt p_i,\wt s_j)$.


Тогда по теореме 1


$$


{\cal F}_{\Theta,Q}


((B^{s_j}_{p_{i,j},1,(1)})_{i=0,\,1,\,j=\const})_{j=0,\,1}=


B^s_{p,q_2,(q_1)}


$$


и


$$


{\cal F}_{\Theta,Q}


((B{}^{\wt s_j}_{\wt p_{i,j},\infty,


(\infty)})_{i=0,\,1,\,j=\const})_{j=0,\,1}=B{}^{\wt s}_{\wt p,q_2,(q_1)}.


$$


Выбирая $q_1=q_2=p$, получим утверждение $T: B^s_{p,p,(p)}


\rightarrow B{}^{\wt s}_{\wt p, p,(p)}$, но в силу неравенства


$p\le \wt p$


$$


B{}^{\wt s}_{\wt p, p,(p)} \subset


B{}^{\wt s_j}_{\wt p, \wt p,(\wt p)}= B^{\wt s}_{\wt p}


\Rightarrow T:B^s_p\rightarrow


B{}^{\wt s}_{\wt p}. \quad\square


$$





Теорема остается верной, если слабые условия типа (3) заменить на


слабые условия одного из следующих типов:


\begin{alignat*}{2}


&\text{a)}&\ T:BL^{s_i,k_j}_{p_j,1} &\rightarrow


BL{}^{\wt s_i,\wt k_j}_{\wt p_j,\infty}, \quad k_j,\ \wt k_j 


\text{ --- любые числа}; \\


%


&\text{б)} &\ T:B^{s_i}_{p_j,1} &\rightarrow


B{}^{\wt s_i}_{\wt p_j,\infty}; \\


%


&\text{в)} &\ T:F^{s_i}_{p_j,1} &\rightarrow


F{}^{\wt s_i}_{\wt p_j,\infty}; \\


%


&\text{г)} &\ T:W^{s_i}_{p_j} &\rightarrow


W{}^{\wt s_i}_{\wt p_j}.


\end{alignat*}





Для любого множества $G\subset R_2$ обозначим через $G^0$


внутренность выпуклой оболочки множества $G$.





\begin{thm}


Пусть в любой точке множества $G\subset(0,\,1)\times R_1$


выполняются слабые условия типа $(3)$. Тогда для


любой внутренней точки $(1/p,s)$ множества $G$ выполняются условия


$$


T:B^s_{p,q_2,(q_1)} \rightarrow B{}^{\wt s}_{\wt p,q_2,(q_1)}.


$$


Если, кроме того, $p\le \wt p$, то


$$


T:B^s_p\rightarrow B{}^{\wt s}_{\wt p}.


$$


\end{thm}





\begin{pf}


Выберем любые точки $(1/p_0,\, s_0),\ (1/p_1,\,s_1) \in G$.


По хорошо известным формулам ([8], теорема 2.4.1) при $p_0=p_1$\;


$(B^{s_0}_{p_0,u}, B^{s_1}_{p_1,u})_{\theta,u} = B^s_{p,u,(u)}$


или $B^s_{p,u}$. Очевидно,


\begin{align*}


(B^{s_0}_{p_0,1,(1)}, B^{s_1}_{p_1,1,(1)})_{\theta,1} &\subset


(B^{s_0}_{p_0,1}, B^{s_1}_{p_1,1})_{\theta,1}, \\


%


(B{}^{\wt s_0}_{\wt p_0,\infty },


B{}^{\wt s_1}_{\wt p_1,\infty })_{\theta,\infty } &\subset


(B{}^{\wt s_0}_{\wt p_0,\infty, (\infty ) },


B{}^{\wt s_1}_{\wt p_1,\infty, (\infty) })_{\theta,\infty }.


\end{align*}


Поэтому из слабых условий в точках $(1/p_i,s_i)$, $i=0,\,1$, следует


$$


T:B^s_{p,1,(1)}\rightarrow B{}^{\wt s}_{\wt p, \infty,(\infty )}


$$


в любой точке $(1/p,s)$ отрезка, соединяющего эти точки, а значит,


слабые условия выполняются в любой точке множества $\ov G^0$.





Пусть теперь $(1/p,s)$ --- произвольная точка из $ \ov G^0$. Для


нее можно указать квадрат со сторонами, параллельными осям координат


так, чтобы его вершины принадлежали $\ov G^0$ и точка $(1/p,s)$


была центром симметрии. Тогда мы получим нужное утверждение с


помощью теоремы 3 при $\Theta = (1/2,1/2)$.


\end{pf}





\section{Билинейные операторы и тензорные произведения}





Декартовы и тензорные произведения банаховых пространств --- еще одна


область, в которой применение многомерных функторов оправдано. Пусть


$(A_0,\,A_1)$ и $(B_0,\,B_1)$ --- интерполяционные пары. Через $W(f,g)$


обозначим билинейный оператор, отображающий декартовы произведения


$A_i\times B_j$ в банаховы пространства $E_{i,j}$, $i=0,\,1$; $j=0,\,1$.


Фиксируя элемент $g\in B_j$, получим линейный оператор


$Tf:=W(f,g)$ ($g=\const$). Очевидно, $T:A_i\rightarrow E_{i,j}$,


следовательно,


$T:A_{\theta, q}\rightarrow (E_{0,j},E_{1,j})_{\theta,q}$.


Аналогичным образом, интерполируя по второму ``аргументу'', получим


интерполяционную теорему.





\begin{thm}


Если билинейный оператор $W$ отображает $A_i\times B_j$


в пространство $E_{i,j}$ при $i=0,\,1;$ $j=0,\,1$, то $W$


отображает $A_{\theta_1,q_1}\times B_{\theta_2,q_2}$ в


$$


{\cal F}_{\Theta,Q}((E_{i,j})_{j=0,\,1})_{i=0,\,1}.


$$


\end{thm}





{\bf Следствие.}


Если линейный оператор $T:A_i\otimes_{\pi}B_j\rightarrow E_{i,j}$, то


$$


T:A_{\theta_1,q_1}\otimes_{\pi} B_{\theta_2,q_2} \rightarrow


{\cal F}_{\Theta,Q}((E_{i,j})_{j=0,\,1})_{i=0,\,1}.


$$


\medskip





Заметим, что при ``одномерной'' интерполяции билинейных операторов


и тензорных произведений встречаются определенные трудности. Пусть


${\cal B\cal L}(A\times B,C)$ обозначает множество непрерывных билинейных


операторов. Для комплексного метода известен простой и исчерпывающий


результат [8],


если $A_0\subset A_1$, $B_0\subset B_1$, $C_0\subset C_1$, то


$ {\cal B\cal L}(A_j\times B_j,C_j) \subset {\cal B\cal L}([A_0,A_1]_{\theta}


\times [B_0,B_1]_{\theta },[C_0,C_1]_{\theta})$.


Для вещественных методов ситуация сложнее.


\medskip





{\bf Пример 1.} Существует билинейный оператор $W$, который


отображает $A_i\times B_i $ в $E_i$, $i=0,\,1$, но не отображает


$A_{\theta, q}\times B_{\theta, q}$ в $E_{\theta, q}$.


\medskip





Пусть $A_0\times B_0 =L_1\times L_{\infty }$,


$A_1\times B_1=L_{\infty }\times L_1$, $E_0=E_1=R$


и оператор $W$ определяется равенством


$$


W(f,g):=\int_{-\infty}^{\infty }f(x)g(x)\,dx.


$$


Очевидно, оператор $W$ отображает $L_1\times L_{\infty }$ и


$L_{\infty }\times L_1$ в $R$, но $W$ не отображает


$(L_1,L_{\infty})_{\theta,q}\times (L_{\infty },L_1)_{\theta, q}


=L_{p,q}\times L_{p',q}$ на $R$ по крайней мере, если $q<2$, т.\,к.


в этом случае $(L_{p,q})^1=L_{p',q'}\not\supset L_{p',q}$.


\medskip





{\bf Пример 2.} Существует линейный оператор $T$, который


отображает $A_i\otimes_{\pi}B_i$, $i=0,\,1$, на $E_i$, но не


отображает $A_{\theta,q}\otimes_{\pi}B_{\theta,q}$ на


$E_{\theta,q}$.


\medskip





Пусть $A_0=L_1(R)$, $A_1=L_{\infty }(R)$,


$B_0=L_{\infty }(R)$, $B_1=L_1(R)$. Пространства


$A_i\otimes_{\pi}B_i$ состоят из функций вида


$f(x,y)= \sum \varphi_j(x) \psi_j(y)$. Определим функционал


$Tf=\int\limits_R f(x,x)\,dx$. Так как из


равенства $\sum \varphi_j(x) \psi_j(y)= \sum \wt \varphi_j(x)


\wt \psi_j(y)$ следует $\sum \varphi_j(x) \psi_j(x)= \sum \wt \varphi_j(x)


\wt \psi_j(x)$, то $Tf$ не зависит от представления функции


$f(x,y)$. Выберем функцию $f(x,y) = \varphi(x) \psi (y)$ так, чтобы


$\varphi \in L_{p,q}$, $\psi \in L_{p',q}$, $q<2$, $\int\limits_R\varphi


(x)\psi (x)\,dx =\infty$. Тогда $T:A_i\otimes_{\pi}B_i


\rightarrow R$, но $Tf=\infty$, хотя


$f\in L_{p,q}\bigotimes_{\pi}L_{p',q}=


A_{\theta,q}\otimes_{\pi}B_{\theta, q}$.





Теорему 5 можно дополнить описанием пространства


${\cal F}_{\Theta,Q}(E_{i,j})$ в случае, когда $E_{i,j}$


образуют одномерное семейство $E_{\theta,q}$.





\begin{thm}


Пусть пространства $E_{0,0}$ и $E_{1,1}$ образуют


интерполяционную пару, а пространства $E_{0,1}$ и $E_{1,0}$


определяются равенствами


$E_{0,1}=(E_{0,0},\;E_{1,1})_{\alpha_0,q_0}$,


$E_{0,1}=(E_{0,0},\;E_{1,1})_{\alpha_1,q_1}$, тогда


${\cal F}_{\Theta,Q} ((E_{i,j})_j)_i =(E_{0,0},\;E_{1,1})_{\beta,\,q_2}$,


где $\beta=(1-\theta_2)\theta_1 \alpha_0+\theta_2(1-\theta_1)\alpha_1$.


\end{thm}





\section{Интегральные операторы в пространствах Бесова}





Рассмотрим интегральный оператор $Kf:= \int\limits_{R_n}


k(x,y)\,f(y)\,dy$. Пусть $E$ --- банахово пространство, элементы которого


суть функции, заданные на $R_n$. Тогда для функции $f(x,y)$,


$x\in R_n$, $y\in R_n$, через $\| f\mid E(x)\|$ обозначим норму от функции


$f$, вычисленную по переменной $x$ при фиксированной переменной $y$. Как


известно ([10], c.\,65),


$\Big\|\int\limits_{R_n} u(x,y)\,dx \mid E\Big\| \le \int\limits_{R_n}


\|u(x,y)\mid E(y)\|\,dx$ для любой банаховой структуры, обладающей свойством


Фату. Можно показать, что это верно также в случае, когда


$E=B^s_{p,q}$.





\begin{lem}


Пусть $1<p<\infty$, $0<s<\infty$,


$\|k(x,y)\mid B^s_{p,\infty}(x)\| \in L_{p',\infty}$.


Тогда интегральный оператор


$$


Kf=\int_{R_n} k(x,y)\, f(y)\, dy


$$


действует из $L_{p,1}$ в $B^s_p$.


\end{lem}





\begin{pf}


Воспользуемся следующей формулой для нормы


пространства Бесова:


$$


\| f\mid B^s_{p,\infty }\| =\| f\mid L_p\| +\sum_{j=1}^m \| h^{-s}


\Delta^{\ell }_{h,j}f \mid L_{\infty }[L_p] \|


$$


при $m>s$,


$\Delta_{h,j}f=f(x_1,\dotsc,x_j+h,\dotsc,x_n)-


f(x_1,\dotsc,x_j,\dotsc,x_n)$,


$\Delta^{\ell}_{h,j}= \Delta_{h,j} (\Delta^{\ell -1}_{h,j})$,


$\ell >s$. Тогда


\begin{multline*}


\bigg\| \int_{R_n} k(x,y)\,f(y)\,dy \mid B^s_{p,\infty } (y)\bigg\| = \\


%


= \bigg\| \int_{R_n} k(x,y)\,f(y)\,dy \mid L_p \bigg\| + \sum^m_{j=1}


\bigg\|h^{-s}\Delta^{\ell }_{h,j} \int_{R_n}


k(x,y)\,f(y)\,dy\mid L_{\infty} [L_{p,\infty }]\bigg\| = \\


%


= \bigg\| \int_{R_n} k(x,y)\,f(y)\,dy \mid L_p \bigg\| + \sum^m_{j=1}


\bigg\|\int_{R_n}h^{-s}\Delta^{\ell }_{h,j} k(x,y)\,f(y)\,dy


\mid L_{\infty }[L_{p,\infty }]\bigg\| \le \\


%


\le \int_{R_n} \| k(x,y)\mid L_p\|\, |f(x)|\,dy + \int_{R_n}


\sum^m_{j=1}\| h^{-s}\Delta^{\ell }_{h,j} k(x,y)\mid


L_{\infty}[L_p]\|\,|f(y)|\,dy = \\


%


= \int_{R_n} \|k(x,y)\mid B^s_{p,\infty }(x)\| \,|f(y)|\, dy \le


\big\| \,\|k(x,y)\mid B^s_{p,\infty }(x)\|\mid


L_{p',\infty}\big\|\, \|f|L_{p,1}\|.


\end{multline*}


Таким образом,


$\|Kf\mid B^s_{p,\infty }\| \le \big\|\,\|k(x,y)\mid B^s_{p,\infty}


(x)\|\mid L_{p',\infty}\big\| \, \|f\mid L_{p,1}\|$.


Следовательно, оператор $K$ действует из пространства $L_{p,1}$ в


пространство Бесова $B^s_{p,\infty }$.


\end{pf}





\begin{lem}


Если $1<p<\infty$, $-\infty <s<\infty$, то из


условия $K:L_{p,1} \rightarrow B^s_{p,\infty }$ следует


$K:B^0_{p,1,(1)} \rightarrow B^s_{p,\infty,(\infty )}$.


\end{lem}





\begin{pf}


Достаточно проверить справедливость вложения


$B^0_{p,1,(1)} \subset L_{p,1}$. Выберем любые $p_0$ и $p_1$ такие,


что $1<p_i<\infty$, $p_0 \ne p_1$ и $1/p=0,5/p_0 +0,5/p_1$. Заметим,


что для любого $q \in (1,\infty)$ справедливы вложения


$\ell^s_1[L_q] \subset L_q[\ell^s_1]\subset L_q[\ell^s_2]$.


Следовательно, $B^0_{q,1} \subset F^0_{q,1} \subset F^0_{q,2} = L_q$.


Применим интерполяционные функторы $(\cdot,\, \cdot )_{1/2,1}$ к


пространствам $B^0_{p_i,1}$ и $L_{p_i}$. Тогда по формуле


(2.4.1/5) из [8]


$$


(B^0_{p_0,1}, B^0_{p_1,1})_{1/2,1}=B^0_{p,1,(1)} \quad \text{и} \quad


(L_{p_0},L_{p_1} )_{1/2,1} =L_{p,1}.


$$


Поэтому из $B^0_{p_i,1} \subset L_{p_i}$ следует


$B^0_{p,1,(1)}\subset L_{p,1}$.


\end{pf}





\begin{thm}


Пусть $K$ --- интегральный оператор с ядром $k(x,y)$,


удовлетворяющим условиям


$$


\| k(x,y)\mid B^{s_i}_{p_i,\infty }\| \in L_{p'_i,\infty }, \quad


i=0,1,2,


$$


где $0<s_i<\infty$, $1<p_i<\infty$ и $(1/p,s)$ --- внутренняя точка


треугольника $A_i(1/p_i,s_i)$. Тогда оператор $K:L_p\rightarrow B^s_p$,


следовательно, $K:B^{\wt s}_p \rightarrow B^s_p$ при любом


$\wt s >0$.


\end{thm}





\begin{pf}


По лемме 4 оператор $K$ удовлетворяет условиям


$K:L_{p_i,1}\rightarrow B^{s_i}_{p_i,\infty }$. Следовательно,


оператор $K$ удовлетворяет слабым условиям


$K: B^0_{p_i,1,(1)}\rightarrow B^{s_i}_{p_i,\infty,(\infty )}$.


Сначала с помощью


``одномерной интерполяции'' покажем, что для любой точки треугольника


$(1/p,s)$ выполняются слабые условия. Для данной точки $(1/p,s)$


выберем такое $\varepsilon >0$, чтобы точки


$E_j(1/p \pm \varepsilon,\ s \pm \varepsilon)$


принадлежали треугольнику $A_0A_1A_2$.


Применим многомерный функтор ${\cal F}_{\Theta,Q}$ к пространствам,


соответствующим точкам $E_j$. Тогда


${\cal F}_{\Theta,Q} (B^{s\pm\varepsilon}_{p_j,\infty})=B^s_p$


при $\Theta = (1/2,1/2)$,


$Q=(p,p)$. В то же время ${\cal F}_{\Theta,Q} (L_{p_i,1})=L_p$. Поэтому


$K:L_p \rightarrow B^s_p$. Из вложений $B^{\wt s}_p \subset L_p$


следует $K:B^{\wt s }_p \rightarrow B^s_p$.


\end{pf}





\begin{thebibliography}{9}





\bibitem{1} 


Sparr G. {\em Interpolation of several Banach spaces} // Ann. Math. 


Pura Appl. -- 1976. -- V.\,99. -- P.\,247--316. 


\bibitem{2} 


Fernandez D.L. {\em Lorentz spaces with mixed norms} // J. Funct. Anal. 


-- 1977. -- V.\,25. -- \No\,2. -- P.\,128--146. 


\bibitem{3} 


Fernandez D.L. {\em Interpolation of $2^n$ Banach spaces} // Studia 


Math. -- 1979. -- V.\,65. -- \No\,2. -- P.\,175--201. 


\bibitem{4} 


Cobos F., Peetre J. {\em Interpolation of compact operators$:$ the 


dimensional case} // Proc. Lond. Math. Soc. -- 1991. -- V.\,63. -- 


P.\,371--400. 


\bibitem{5} 


Milman M {\em On interpolation of $2^n$ Banach spaces and Lorentz 


spaces with Mixed norms} // J. Funct. Anal. -- 1981. -- V.\,41. -- 


P.\,1--7. 


\bibitem{6} 


Асекритова И.У. {\em Вещественный метод интерполяции для конечных 


наборов банаховых пространств} // Исследов. по теории функций многих 


веществ. переменных. -- Ярославль, 1981. -- С.\,9--17. 


\bibitem{7} 


Cwikel M., Janson S. {\em Real and complex interpolation for finite and 


infinite families of $B$-spaces} // Advances in Matem. -- 1987. -- 


V.\,66. -- P.\,234--290. 


\bibitem{8} 


Трибель Х. {\em Теория интерполяции, функциональные пространства, 


дифференциальные операторы}. -- М.: Мир, 1980. -- 664 с. 


\bibitem{9} 


Крепкогорский В.Л. {\em Интерполяция в пространствах Лизоркина--Трибеля 


и Бесова} // Матем. сб. -- 1994. -- Т.\,185. -- \No\,7. -- С.\,63--76. 





\end{thebibliography}





\vskip 2cm





\postup{\it Казанский филиал военного}{$\qquad$\it Поступили} 


\postup{\it артиллерийского университета}{{\it первый вариант} 24.06.1997} 


\postup{}{{\it окончательный вариант} 18.06.1999} 





\label{end}


\end{document}


